\documentclass[a4paper]{article}
\usepackage[fontset=fandol]{ctex}
\usepackage[margin=1.65cm]{geometry}
\usepackage{amsmath,amssymb}
\usepackage{graphicx}
\usepackage{hyperref}
\usepackage{longtable,booktabs}
\usepackage{xcolor}
\usepackage{enumitem}
\usepackage{listings}
\lstset{basicstyle=\ttfamily\small,breaklines=true,columns=fullflexible,keepspaces=true}
\hypersetup{colorlinks=true,linkcolor=blue,urlcolor=blue}
\setlist[itemize]{leftmargin=1.25em,itemsep=0.18em}
\setlength{\parskip}{0.45em}
\setlength{\parindent}{2em}
\sloppy
\title{第 02 讲：PyTorch、einops 与资源核算}
\author{CS336 Spring 2026 public videos and official course materials}
\date{}
\begin{document}
\maketitle
\tableofcontents
\newpage
\section{本章导读}
本章讨论 \textbf{PyTorch、einops 与资源核算}。CS336 的第一层目标不是把现成大模型当作黑箱使用，而是把 language model 从输入文本、训练目标、模型结构、数据、系统和评估这些部件重新拆开。只有拆开之后，读者才能判断一个设计选择究竟改变了什么。
设想你看到一个训练脚本声称能在若干 H100 上训练 70B 模型。判断它是否可信，第一步不是读模型名字，而是问：参数、梯度、optimizer state 和 activation 各占多少显存？每一步大约多少 FLOPs？瓶颈是 tensor core 还是 HBM bandwidth？
本章对应 CS336 Spring 2026 第 02 个公开视频，并结合课程页提供的可解析资料。视频和课程页给出事实边界；正文把这些材料改写为可自学的教材章节。
阅读本章时，先理解问题为什么存在，再看定义、公式和实现形态，随后通过例题和练习检查自己是否能独立复现这条 reasoning chain。
\begin{figure}[h]
\centering
\includegraphics[width=0.58\linewidth]{source_anchor.jpg}
\caption{第 02 讲的课程图像锚点。}
\end{figure}
\subsection{本章知识路线}
\begin{itemize}
\item 本章首先说明 PyTorch、einops 与资源核算 为什么是 language modeling pipeline 中不可跳过的一环。
\item 其次定义本章反复使用的术语，并说明这些术语对应的计算对象或系统对象。
\item 随后用公式和伪代码表达主要机制，让概念能够落到可计算的变量上。
\item 最后给出例题、风险讨论和练习，便于读者检查自己是否真正掌握了机制。
\end{itemize}
\subsection{结构图}
下面的图用于帮助读者定位本章的核心结构。先看概念之间的依赖，再回到正文中的公式、伪代码和例题。
\begin{figure}[h]
\centering
\includegraphics[width=0.72\linewidth]{figures/source_grounded_concept_map.png}
\caption{第 02 讲概念图：本章问题、术语和机制的关系。}
\end{figure}
\begin{figure}[h]
\centering
\includegraphics[width=0.72\linewidth]{figures/source_grounded_checklist.png}
\caption{第 02 讲学习路径：从问题定义进入公式、实现、实验和练习。}
\end{figure}
\section{本章主线}
本章可以沿着三条线索阅读：课程为什么把这个主题放在这里，它改变了哪些变量，以及怎样用小实验检查这些变量。
\subsection{本章要解决的核心问题}
本讲把 PyTorch tensor 看作所有训练对象的统一表示：参数、梯度、optimizer state、activation 和数据 batch 都是带 shape、dtype、device 的张量。
这一点对应的学习动作是：找出对象，写出变量，说明变量如何被测量。若输入规模、硬件、数据分布或评估协议改变，结论也必须重新检查。
因此，PyTorch、einops 与资源核算 不应被读成若干名词的集合，而应被读成一组可检验的机制。
\subsection{从机制到资源账本}
资源核算不是附属技能。FLOPs 决定训练时间下界，显存决定模型和 batch 是否放得下，内存带宽决定很多 elementwise 或 decode workload 是否会被 HBM 限制。
这一点对应的学习动作是：找出对象，写出变量，说明变量如何被测量。若输入规模、硬件、数据分布或评估协议改变，结论也必须重新检查。
因此，PyTorch、einops 与资源核算 不应被读成若干名词的集合，而应被读成一组可检验的机制。
\subsection{从课堂讲解到可复现实验}
einops 的价值不是语法糖，而是让维度语义显式化。把 batch、sequence、head、hidden 等维度写清楚，可以减少实现错误，也方便做 FLOPs/bytes 的手算。
这一点对应的学习动作是：找出对象，写出变量，说明变量如何被测量。若输入规模、硬件、数据分布或评估协议改变，结论也必须重新检查。
因此，PyTorch、einops 与资源核算 不应被读成若干名词的集合，而应被读成一组可检验的机制。
\section{术语、符号与问题边界}
本章的术语横跨模型、数据和系统。读者需要先知道每个词指向的对象：有些词描述数学目标，有些词描述张量形状，有些词描述硬件瓶颈，还有一些词描述评估协议。术语边界不清，后面的公式和实现就会变成空洞的缩写。
\begin{longtable}{p{0.08\linewidth}p{0.42\linewidth}p{0.42\linewidth}}
\toprule
\textbf{\#} & \textbf{术语} & \textbf{阅读要求}\\
\midrule
1 & tensor（张量） & 需要能说明它指向的对象、参与的公式或代码位置，以及一个典型失败情形。\\
2 & mixed precision（混合精度） & 需要能说明它指向的对象、参与的公式或代码位置，以及一个典型失败情形。\\
3 & FLOPs（浮点运算次数） & 需要能说明它指向的对象、参与的公式或代码位置，以及一个典型失败情形。\\
4 & memory bandwidth（内存带宽） & 需要能说明它指向的对象、参与的公式或代码位置，以及一个典型失败情形。\\
5 & arithmetic intensity（算术强度） & 需要能说明它指向的对象、参与的公式或代码位置，以及一个典型失败情形。\\
6 & roofline analysis（屋顶线分析） & 需要能说明它指向的对象、参与的公式或代码位置，以及一个典型失败情形。\\
\bottomrule
\end{longtable}
这些术语之间有依赖关系。例如 tokenization 改变 sequence length，sequence length 又改变 attention FLOPs、KV cache 和 evaluation token budget；parallelism 改变显存占用，也改变通信瓶颈和故障恢复方式。
\subsection{关键词}
下面的关键词在课程材料中反复出现。它们不代替定义，只提示读者本章哪些对象需要在后文持续跟踪。
\begin{itemize}
\item function, memory, flops, compute, training, intensity, parameters, activations, data, accelerator, pass, arithmetic, memory-bound, fp32
\end{itemize}
\section{张量、dtype 与显存账本}
现在进入 张量、dtype 与显存账本。本节关心的不是名称本身，而是它把 PyTorch、einops 与资源核算 中的哪些对象变成可计算、可测量或可优化的量。
\subsection{为什么要讨论 张量、dtype 与显存账本}
课程材料用 rank、shape、dtype、device 解释 PyTorch tensor；参数、梯度、optimizer state、activation 和 data 都是 tensor。
fp32、fp16、bf16、fp8、fp4 的差别首先体现在每个元素的字节数和动态范围，而不是抽象精度标签。
显存估算的基本动作是：元素个数乘以 element size，再乘以副本数量。AdamW 通常还要记参数、梯度和两个 fp32 optimizer states。
张量、dtype 与显存账本 的作用是把 PyTorch、einops 与资源核算 中的一个抽象问题变成可计算对象：哪些变量进入公式，哪些状态进入代码，哪些条件会让结果失效。
\subsection{张量、dtype 与显存账本 的机制}
课程材料用 rank、shape、dtype、device 解释 PyTorch tensor；参数、梯度、optimizer state、activation 和 data 都是 tensor。 这给出本节的对象和边界。读者需要把其中的名词对应到张量、数据记录、训练状态或系统状态，而不是停留在缩写层面。
fp32、fp16、bf16、fp8、fp4 的差别首先体现在每个元素的字节数和动态范围，而不是抽象精度标签。 因而同一个方法在小规模示例、单机实验和集群训练中的意义并不相同。比较方法时，问题规模、硬件条件、数据分布和评估口径必须同时给出。
显存估算的基本动作是：元素个数乘以 element size，再乘以副本数量。AdamW 通常还要记参数、梯度和两个 fp32 optimizer states。 这使本节具有可检验性：如果一个结论不能在公式、伪代码、profile trace、benchmark 或 ablation 中表现出来，它就还只是直觉。
因此，张量、dtype 与显存账本 应当被理解为一组可以实现和测试的假设。CS336 反复强调 from scratch，正是因为只有能被复现的机制，才可能被稳定改进。
\subsection{张量、dtype 与显存账本 的数学表达}
\[\text{bytes}=\#\text{elements}\times\text{bytes/element}\times\#\text{copies}\]
符号说明：copies 可以来自参数、梯度、momentum、variance、activation checkpoint 或并行副本；它是系统设计中最容易漏算的一项。
这个关系式给出了本节最小的数学骨架。读者应检查每个变量的单位、取值范围和可观测性；如果某个量只能间接估计，后续实验就必须说明 proxy 是什么。
例如，validation loss 常被用作模型质量的 proxy，tokens/s 用作吞吐 proxy，Jaccard 或 MinHash 用作近重复 proxy，reward model score 用作偏好 proxy。proxy 不是目标本身；它只在评估协议清楚时才可引用。
\subsection{从 张量、dtype 与显存账本 到程序}
本节机制可以写成下面的伪代码。它保留课程材料中的计算顺序和状态变化，但省略了工程代码中的错误处理、并行细节和框架样板。
\begin{lstlisting}[language=Python]
def tensor_bytes(shape, bytes_per_element, copies=1):
    return prod(shape) * bytes_per_element * copies
\end{lstlisting}
读这段伪代码时，重点不是语法，而是状态如何变化：输入是什么，中间变量是什么，主要计算或数据移动发生在哪里，哪些边界条件会破坏结果。
\subsection{边界条件与常见误解}
\begin{itemize}
\item activation memory 依赖 batch size、sequence length 和 checkpointing 策略，不能只看参数量。
\item 低精度 dtype 节省显存和带宽，但 dynamic range 与 accumulation 精度会影响训练稳定性。
\item 不要把小实验中的局部结论自动外推到 frontier-scale；scale、数据分布和硬件拓扑都会改变结论。
\item 不要把 benchmark 或 profile 的单个数字当作机制解释；必须同时记录实验设置、输入形状、版本和评估口径。
\end{itemize}
\subsection{与本章其他概念的关系}
张量、dtype 与显存账本 与本章其他部分的关系是互补的：术语表给出对象名称，公式给出最小数学结构，伪代码给出执行顺序，例题则把这些内容压缩成一个可检查的计算。
\subsection{怎样检验 张量、dtype 与显存账本}
检验 张量、dtype 与显存账本 的第一步是确定测量对象。这里的测量对象通常不是一个抽象名词，而是与 tensor、dtype、fp32、bf16、optimizer 有关的张量、样本、请求、奖励、通信操作或评估记录。如果对象没有被具体化，后面的实验就只能得到描述性观察，不能得到可复现结论。
第二步是构造最小输入。最小输入应该足够小，使读者能手算或打印关键中间量；同时又要保留本节机制。对于 张量、dtype 与显存账本，最小输入不追求逼真规模，而追求暴露因果关系：改变一个变量时，哪些输出、成本或错误会随之改变。
第三步是写出资源账本。账本可以是 FLOPs、bytes moved、显存峰值、通信量、token 数、样本数、reward 分布或 benchmark 分数。符号说明：copies 可以来自参数、梯度、momentum、variance、activation checkpoint 或并行副本；它是系统设计中最容易漏算的一项。 这类说明应被转化为单位明确的数量关系；如果数量只能间接观测，就必须说明使用了什么 proxy。
第四步是设置对照。一个好对照只改变一个因素，例如 batch size、sequence length、vocabulary size、attention heads、filter threshold、learning rate、decoding 参数或 verifier。如果多个因素同时改变，实验最多说明系统表现不同，不能说明是哪一个机制在起作用。
第五步是观察失败条件。教材中的成功路径往往最容易记住，但工程中真正决定方法边界的是失败案例。本节至少要记住这些限制：activation memory 依赖 batch size、sequence length 和 checkpointing 策略，不能只看参数量；低精度 dtype 节省显存和带宽，但 dynamic range 与 accumulation 精度会影响训练稳定性。复习时可以把这些限制改写成反例测试。
最后，把实验结论放回 PyTorch、einops 与资源核算 的整体结构中。张量、dtype 与显存账本 不是孤立技巧：它会影响训练目标、数据选择、系统成本或评估解释中的至少一项。能说清楚这种影响，才说明读者已经从名词记忆进入机制理解。
\subsection{常见混淆：张量、dtype 与显存账本}
tensor（张量）是参数、激活、梯度和数据的共同表示。在“张量、dtype 与显存账本”中，这个概念用于解释 PyTorch、einops 与资源核算 的一个具体机制。
dtype 在本节中应被看作一个技术对象，而不是一个孤立名词。需要判断它指向模型、数据、系统还是评估，并说明它在“张量、dtype 与显存账本”中改变了哪个变量或约束。
fp32 在本节中应被看作一个技术对象，而不是一个孤立名词。需要判断它指向模型、数据、系统还是评估，并说明它在“张量、dtype 与显存账本”中改变了哪个变量或约束。
bf16 在本节中应被看作一个技术对象，而不是一个孤立名词。需要判断它指向模型、数据、系统还是评估，并说明它在“张量、dtype 与显存账本”中改变了哪个变量或约束。
optimizer 在本节中应被看作一个技术对象，而不是一个孤立名词。需要判断它指向模型、数据、系统还是评估，并说明它在“张量、dtype 与显存账本”中改变了哪个变量或约束。
activation 在本节中应被看作一个技术对象，而不是一个孤立名词。需要判断它指向模型、数据、系统还是评估，并说明它在“张量、dtype 与显存账本”中改变了哪个变量或约束。
这些概念的共同点是：只有进入公式、代码或评估协议之后才有明确含义。单独背诵名称不会帮助判断一个训练 run、推理系统或数据 pipeline 是否正确。
\subsection{张量、dtype 与显存账本 的小结}
\begin{itemize}
\item 张量、dtype 与显存账本 的任何性能结论都必须同时报告问题规模和评估口径，否则不可比较。
\item 如果一个方法声称降低主要成本，它通常会把成本转移到另一个变量，例如显存、通信、数据质量、训练稳定性或评估复杂度。
\item 公式和伪代码提供推理骨架；真实系统仍需要 profiler、ablation 和 held-out evaluation。
\end{itemize}
\section{FLOPs 与训练时间的 napkin math}
现在进入 FLOPs 与训练时间的 napkin math。本节关心的不是名称本身，而是它把 PyTorch、einops 与资源核算 中的哪些对象变成可计算、可测量或可优化的量。
\subsection{为什么要讨论 FLOPs 与训练时间的 napkin math}
视频用 70B 参数、15T tokens、1024 H100 这类例子展示快速估算训练时间的方式。
经验公式 \textbackslash{}(6ND\textbackslash{}) 把 forward、backward 和 gradient computation 合并成每 token 每参数约 6 FLOPs 的估计。
模型 FLOPs 利用率（MFU）把理论峰值和真实训练吞吐连接起来，是判断 kernel、parallelism 和 data pipeline 是否有效的重要指标。
FLOPs 与训练时间的 napkin math 的作用是把 PyTorch、einops 与资源核算 中的一个抽象问题变成可计算对象：哪些变量进入公式，哪些状态进入代码，哪些条件会让结果失效。
\subsection{FLOPs 与训练时间的 napkin math 的机制}
视频用 70B 参数、15T tokens、1024 H100 这类例子展示快速估算训练时间的方式。 这给出本节的对象和边界。读者需要把其中的名词对应到张量、数据记录、训练状态或系统状态，而不是停留在缩写层面。
经验公式 \textbackslash{}(6ND\textbackslash{}) 把 forward、backward 和 gradient computation 合并成每 token 每参数约 6 FLOPs 的估计。 因而同一个方法在小规模示例、单机实验和集群训练中的意义并不相同。比较方法时，问题规模、硬件条件、数据分布和评估口径必须同时给出。
模型 FLOPs 利用率（MFU）把理论峰值和真实训练吞吐连接起来，是判断 kernel、parallelism 和 data pipeline 是否有效的重要指标。 这使本节具有可检验性：如果一个结论不能在公式、伪代码、profile trace、benchmark 或 ablation 中表现出来，它就还只是直觉。
因此，FLOPs 与训练时间的 napkin math 应当被理解为一组可以实现和测试的假设。CS336 反复强调 from scratch，正是因为只有能被复现的机制，才可能被稳定改进。
\subsection{FLOPs 与训练时间的 napkin math 的数学表达}
\[\text{training FLOPs}\approx 6ND,\qquad \text{days}\approx \frac{6ND}{F_{\text{gpu}}\cdot n_{\text{gpu}}\cdot \text{MFU}\cdot 86400}\]
符号说明：\(N\) 是参数量，\(D\) 是训练 token 数，\(F_{\text{gpu}}\) 是单 GPU 理论 FLOPs/s，MFU 是有效利用率。
这个关系式给出了本节最小的数学骨架。读者应检查每个变量的单位、取值范围和可观测性；如果某个量只能间接估计，后续实验就必须说明 proxy 是什么。
例如，validation loss 常被用作模型质量的 proxy，tokens/s 用作吞吐 proxy，Jaccard 或 MinHash 用作近重复 proxy，reward model score 用作偏好 proxy。proxy 不是目标本身；它只在评估协议清楚时才可引用。
\subsection{从 FLOPs 与训练时间的 napkin math 到程序}
本节机制可以写成下面的伪代码。它保留课程材料中的计算顺序和状态变化，但省略了工程代码中的错误处理、并行细节和框架样板。
\begin{lstlisting}[language=Python]
total_flops = 6 * num_parameters * num_tokens
seconds = total_flops / (gpu_flops * num_gpus * mfu)
\end{lstlisting}
读这段伪代码时，重点不是语法，而是状态如何变化：输入是什么，中间变量是什么，主要计算或数据移动发生在哪里，哪些边界条件会破坏结果。
\subsection{边界条件与常见误解}
\begin{itemize}
\item 这个估算忽略 optimizer、embedding、attention/MLP 比例、通信和 data loading，但足以暴露数量级错误。
\item 不要把 peak FLOPs 当作可持续吞吐；实际 MFU 需要测量。
\item 不要把小实验中的局部结论自动外推到 frontier-scale；scale、数据分布和硬件拓扑都会改变结论。
\item 不要把 benchmark 或 profile 的单个数字当作机制解释；必须同时记录实验设置、输入形状、版本和评估口径。
\end{itemize}
\subsection{与本章其他概念的关系}
FLOPs 与训练时间的 napkin math 与本章其他部分的关系是互补的：术语表给出对象名称，公式给出最小数学结构，伪代码给出执行顺序，例题则把这些内容压缩成一个可检查的计算。
\subsection{怎样检验 FLOPs 与训练时间的 napkin math}
检验 FLOPs 与训练时间的 napkin math 的第一步是确定测量对象。这里的测量对象通常不是一个抽象名词，而是与 FLOPs、MFU、H100、70B、tokens 有关的张量、样本、请求、奖励、通信操作或评估记录。如果对象没有被具体化，后面的实验就只能得到描述性观察，不能得到可复现结论。
第二步是构造最小输入。最小输入应该足够小，使读者能手算或打印关键中间量；同时又要保留本节机制。对于 FLOPs 与训练时间的 napkin math，最小输入不追求逼真规模，而追求暴露因果关系：改变一个变量时，哪些输出、成本或错误会随之改变。
第三步是写出资源账本。账本可以是 FLOPs、bytes moved、显存峰值、通信量、token 数、样本数、reward 分布或 benchmark 分数。符号说明：\textbackslash{}(N\textbackslash{}) 是参数量，\textbackslash{}(D\textbackslash{}) 是训练 token 数，\textbackslash{}(F\_\{ ext\{gpu\}\}\textbackslash{}) 是单 GPU 理论 FLOPs/s，MFU 是有效利用率。 这类说明应被转化为单位明确的数量关系；如果数量只能间接观测，就必须说明使用了什么 proxy。
第四步是设置对照。一个好对照只改变一个因素，例如 batch size、sequence length、vocabulary size、attention heads、filter threshold、learning rate、decoding 参数或 verifier。如果多个因素同时改变，实验最多说明系统表现不同，不能说明是哪一个机制在起作用。
第五步是观察失败条件。教材中的成功路径往往最容易记住，但工程中真正决定方法边界的是失败案例。本节至少要记住这些限制：这个估算忽略 optimizer、embedding、attention/MLP 比例、通信和 data loading，但足以暴露数量级错误；不要把 peak FLOPs 当作可持续吞吐；实际 MFU 需要测量。复习时可以把这些限制改写成反例测试。
最后，把实验结论放回 PyTorch、einops 与资源核算 的整体结构中。FLOPs 与训练时间的 napkin math 不是孤立技巧：它会影响训练目标、数据选择、系统成本或评估解释中的至少一项。能说清楚这种影响，才说明读者已经从名词记忆进入机制理解。
\subsection{常见混淆：FLOPs 与训练时间的 napkin math}
FLOPs（floating-point operations）是训练和推理成本账本的核心单位。在“FLOPs 与训练时间的 napkin math”中，这个概念用于解释 PyTorch、einops 与资源核算 的一个具体机制。
MFU 在本节中应被看作一个技术对象，而不是一个孤立名词。需要判断它指向模型、数据、系统还是评估，并说明它在“FLOPs 与训练时间的 napkin math”中改变了哪个变量或约束。
H100 在本节中应被看作一个技术对象，而不是一个孤立名词。需要判断它指向模型、数据、系统还是评估，并说明它在“FLOPs 与训练时间的 napkin math”中改变了哪个变量或约束。
70B 在本节中应被看作一个技术对象，而不是一个孤立名词。需要判断它指向模型、数据、系统还是评估，并说明它在“FLOPs 与训练时间的 napkin math”中改变了哪个变量或约束。
tokens 在本节中应被看作一个技术对象，而不是一个孤立名词。需要判断它指向模型、数据、系统还是评估，并说明它在“FLOPs 与训练时间的 napkin math”中改变了哪个变量或约束。
这些概念的共同点是：只有进入公式、代码或评估协议之后才有明确含义。单独背诵名称不会帮助判断一个训练 run、推理系统或数据 pipeline 是否正确。
\subsection{FLOPs 与训练时间的 napkin math 的小结}
\begin{itemize}
\item FLOPs 与训练时间的 napkin math 的任何性能结论都必须同时报告问题规模和评估口径，否则不可比较。
\item 如果一个方法声称降低主要成本，它通常会把成本转移到另一个变量，例如显存、通信、数据质量、训练稳定性或评估复杂度。
\item 公式和伪代码提供推理骨架；真实系统仍需要 profiler、ablation 和 held-out evaluation。
\end{itemize}
\section{einops、算术强度与 bottleneck 判断}
现在进入 einops、算术强度与 bottleneck 判断。本节关心的不是名称本身，而是它把 PyTorch、einops 与资源核算 中的哪些对象变成可计算、可测量或可优化的量。
\subsection{为什么要讨论 einops、算术强度与 bottleneck 判断}
einops 通过命名维度减少 PyTorch 代码里 transpose、view、reshape 的认知负担。
arithmetic intensity 衡量每搬运 1 byte 做多少 FLOPs；矩阵乘法通常 compute-bound，elementwise ops 常常 memory-bound。
roofline analysis 给出性能上限：低算术强度受带宽限制，高算术强度受峰值计算限制。
einops、算术强度与 bottleneck 判断 的作用是把 PyTorch、einops 与资源核算 中的一个抽象问题变成可计算对象：哪些变量进入公式，哪些状态进入代码，哪些条件会让结果失效。
\subsection{einops、算术强度与 bottleneck 判断 的机制}
einops 通过命名维度减少 PyTorch 代码里 transpose、view、reshape 的认知负担。 这给出本节的对象和边界。读者需要把其中的名词对应到张量、数据记录、训练状态或系统状态，而不是停留在缩写层面。
arithmetic intensity 衡量每搬运 1 byte 做多少 FLOPs；矩阵乘法通常 compute-bound，elementwise ops 常常 memory-bound。 因而同一个方法在小规模示例、单机实验和集群训练中的意义并不相同。比较方法时，问题规模、硬件条件、数据分布和评估口径必须同时给出。
roofline analysis 给出性能上限：低算术强度受带宽限制，高算术强度受峰值计算限制。 这使本节具有可检验性：如果一个结论不能在公式、伪代码、profile trace、benchmark 或 ablation 中表现出来，它就还只是直觉。
因此，einops、算术强度与 bottleneck 判断 应当被理解为一组可以实现和测试的假设。CS336 反复强调 from scratch，正是因为只有能被复现的机制，才可能被稳定改进。
\subsection{einops、算术强度与 bottleneck 判断 的数学表达}
\[\text{arithmetic intensity}=\frac{\text{FLOPs}}{\text{bytes moved}},\qquad P\le \min(P_{\max},B_{\max}\cdot I)\]
符号说明：\(I\) 是算术强度，\(P_{\max}\) 是峰值计算，\(B_{\max}\) 是内存带宽，\(P\) 是可达到性能。
这个关系式给出了本节最小的数学骨架。读者应检查每个变量的单位、取值范围和可观测性；如果某个量只能间接估计，后续实验就必须说明 proxy 是什么。
例如，validation loss 常被用作模型质量的 proxy，tokens/s 用作吞吐 proxy，Jaccard 或 MinHash 用作近重复 proxy，reward model score 用作偏好 proxy。proxy 不是目标本身；它只在评估协议清楚时才可引用。
\subsection{从 einops、算术强度与 bottleneck 判断 到程序}
本节机制可以写成下面的伪代码。它保留课程材料中的计算顺序和状态变化，但省略了工程代码中的错误处理、并行细节和框架样板。
\begin{lstlisting}[language=Python]
scores = einsum(x, y, 'batch seq hidden, batch cand hidden -> batch seq cand')
if flops / bytes_moved < roofline_knee:
    optimize_memory_traffic()
else:
    optimize_compute_utilization()
\end{lstlisting}
读这段伪代码时，重点不是语法，而是状态如何变化：输入是什么，中间变量是什么，主要计算或数据移动发生在哪里，哪些边界条件会破坏结果。
\subsection{边界条件与常见误解}
\begin{itemize}
\item roofline 是上界模型，不会告诉你具体哪个 kernel 写错了。
\item 维度语义清晰不等于性能最佳；layout 和 contiguous memory 仍然重要。
\item 不要把小实验中的局部结论自动外推到 frontier-scale；scale、数据分布和硬件拓扑都会改变结论。
\item 不要把 benchmark 或 profile 的单个数字当作机制解释；必须同时记录实验设置、输入形状、版本和评估口径。
\end{itemize}
\subsection{与本章其他概念的关系}
einops、算术强度与 bottleneck 判断 与本章其他部分的关系是互补的：术语表给出对象名称，公式给出最小数学结构，伪代码给出执行顺序，例题则把这些内容压缩成一个可检查的计算。
\subsection{怎样检验 einops、算术强度与 bottleneck 判断}
检验 einops、算术强度与 bottleneck 判断 的第一步是确定测量对象。这里的测量对象通常不是一个抽象名词，而是与 einops、einsum、arithmetic intensity、roofline、memory-bound 有关的张量、样本、请求、奖励、通信操作或评估记录。如果对象没有被具体化，后面的实验就只能得到描述性观察，不能得到可复现结论。
第二步是构造最小输入。最小输入应该足够小，使读者能手算或打印关键中间量；同时又要保留本节机制。对于 einops、算术强度与 bottleneck 判断，最小输入不追求逼真规模，而追求暴露因果关系：改变一个变量时，哪些输出、成本或错误会随之改变。
第三步是写出资源账本。账本可以是 FLOPs、bytes moved、显存峰值、通信量、token 数、样本数、reward 分布或 benchmark 分数。符号说明：\textbackslash{}(I\textbackslash{}) 是算术强度，\textbackslash{}(P\_\{\textbackslash{}max\}\textbackslash{}) 是峰值计算，\textbackslash{}(B\_\{\textbackslash{}max\}\textbackslash{}) 是内存带宽，\textbackslash{}(P\textbackslash{}) 是可达到性能。 这类说明应被转化为单位明确的数量关系；如果数量只能间接观测，就必须说明使用了什么 proxy。
第四步是设置对照。一个好对照只改变一个因素，例如 batch size、sequence length、vocabulary size、attention heads、filter threshold、learning rate、decoding 参数或 verifier。如果多个因素同时改变，实验最多说明系统表现不同，不能说明是哪一个机制在起作用。
第五步是观察失败条件。教材中的成功路径往往最容易记住，但工程中真正决定方法边界的是失败案例。本节至少要记住这些限制：roofline 是上界模型，不会告诉你具体哪个 kernel 写错了；维度语义清晰不等于性能最佳；layout 和 contiguous memory 仍然重要。复习时可以把这些限制改写成反例测试。
最后，把实验结论放回 PyTorch、einops 与资源核算 的整体结构中。einops、算术强度与 bottleneck 判断 不是孤立技巧：它会影响训练目标、数据选择、系统成本或评估解释中的至少一项。能说清楚这种影响，才说明读者已经从名词记忆进入机制理解。
\subsection{常见混淆：einops、算术强度与 bottleneck 判断}
einops 在本节中应被看作一个技术对象，而不是一个孤立名词。需要判断它指向模型、数据、系统还是评估，并说明它在“einops、算术强度与 bottleneck 判断”中改变了哪个变量或约束。
einsum 在本节中应被看作一个技术对象，而不是一个孤立名词。需要判断它指向模型、数据、系统还是评估，并说明它在“einops、算术强度与 bottleneck 判断”中改变了哪个变量或约束。
arithmetic intensity 在本节中应被看作一个技术对象，而不是一个孤立名词。需要判断它指向模型、数据、系统还是评估，并说明它在“einops、算术强度与 bottleneck 判断”中改变了哪个变量或约束。
roofline 在本节中应被看作一个技术对象，而不是一个孤立名词。需要判断它指向模型、数据、系统还是评估，并说明它在“einops、算术强度与 bottleneck 判断”中改变了哪个变量或约束。
memory（内存/显存）约束决定 batch、sequence length、optimizer state 和 KV cache 能否放下。在“einops、算术强度与 bottleneck 判断”中，这个概念用于解释 PyTorch、einops 与资源核算 的一个具体机制。
这些概念的共同点是：只有进入公式、代码或评估协议之后才有明确含义。单独背诵名称不会帮助判断一个训练 run、推理系统或数据 pipeline 是否正确。
\subsection{einops、算术强度与 bottleneck 判断 的小结}
\begin{itemize}
\item einops、算术强度与 bottleneck 判断 的任何性能结论都必须同时报告问题规模和评估口径，否则不可比较。
\item 如果一个方法声称降低主要成本，它通常会把成本转移到另一个变量，例如显存、通信、数据质量、训练稳定性或评估复杂度。
\item 公式和伪代码提供推理骨架；真实系统仍需要 profiler、ablation 和 held-out evaluation。
\end{itemize}
\section{综合讨论：从局部机制到完整系统}
本章的主题是 PyTorch、einops 与资源核算，但它不应被读成几个相邻知识点的拼接。张量、dtype 与显存账本、FLOPs 与训练时间的 napkin math、einops、算术强度与 bottleneck 判断 共同回答同一个问题：语言模型系统中哪些对象可以被定义，哪些成本可以被估算，哪些结论必须通过实验或评估来确认。
这些对象包括 tensor（张量）、mixed precision（混合精度）、FLOPs（浮点运算次数）、memory bandwidth（内存带宽）、arithmetic intensity（算术强度）、roofline analysis（屋顶线分析）。读者可以把它们看成一组接口：有的接口连接文本和 token，有的连接模型结构和张量计算，有的连接数据样本和训练目标，有的连接离线评估和在线服务。接口一旦改变，后面的成本、错误类型和优化空间也会随之改变。
\subsection{张量、dtype 与显存账本 在系统中的位置}
张量、dtype 与显存账本 首先提供一个局部解释：课程材料用 rank、shape、dtype、device 解释 PyTorch tensor；参数、梯度、optimizer state、activation 和 data 都是 tensor。 这类解释的价值在于把直觉压缩成可操作对象。一旦对象明确，读者就可以追问它的输入、输出、中间状态和失败条件，而不是只记住方法名。
其次，张量、dtype 与显存账本 会改变至少一个资源或评估变量。它可能改变 sequence length、activation size、memory bandwidth、communication volume、data mixture、reward distribution 或 benchmark protocol。这些变量在小例子里可以手算，在真实训练和服务系统里则需要 profiling、logging 和 held-out evaluation。
最后，张量、dtype 与显存账本 不能脱离边界条件理解。一个关键 caveat 是：activation memory 依赖 batch size、sequence length 和 checkpointing 策略，不能只看参数量。。把 caveat 写出来不是为了削弱结论，而是为了说明结论在哪些输入、硬件、数据和评估条件下才成立。
\subsection{FLOPs 与训练时间的 napkin math 在系统中的位置}
FLOPs 与训练时间的 napkin math 首先提供一个局部解释：视频用 70B 参数、15T tokens、1024 H100 这类例子展示快速估算训练时间的方式。 这类解释的价值在于把直觉压缩成可操作对象。一旦对象明确，读者就可以追问它的输入、输出、中间状态和失败条件，而不是只记住方法名。
其次，FLOPs 与训练时间的 napkin math 会改变至少一个资源或评估变量。它可能改变 sequence length、activation size、memory bandwidth、communication volume、data mixture、reward distribution 或 benchmark protocol。这些变量在小例子里可以手算，在真实训练和服务系统里则需要 profiling、logging 和 held-out evaluation。
最后，FLOPs 与训练时间的 napkin math 不能脱离边界条件理解。一个关键 caveat 是：这个估算忽略 optimizer、embedding、attention/MLP 比例、通信和 data loading，但足以暴露数量级错误。。把 caveat 写出来不是为了削弱结论，而是为了说明结论在哪些输入、硬件、数据和评估条件下才成立。
\subsection{einops、算术强度与 bottleneck 判断 在系统中的位置}
einops、算术强度与 bottleneck 判断 首先提供一个局部解释：einops 通过命名维度减少 PyTorch 代码里 transpose、view、reshape 的认知负担。 这类解释的价值在于把直觉压缩成可操作对象。一旦对象明确，读者就可以追问它的输入、输出、中间状态和失败条件，而不是只记住方法名。
其次，einops、算术强度与 bottleneck 判断 会改变至少一个资源或评估变量。它可能改变 sequence length、activation size、memory bandwidth、communication volume、data mixture、reward distribution 或 benchmark protocol。这些变量在小例子里可以手算，在真实训练和服务系统里则需要 profiling、logging 和 held-out evaluation。
最后，einops、算术强度与 bottleneck 判断 不能脱离边界条件理解。一个关键 caveat 是：roofline 是上界模型，不会告诉你具体哪个 kernel 写错了。。把 caveat 写出来不是为了削弱结论，而是为了说明结论在哪些输入、硬件、数据和评估条件下才成立。
\subsection{把本章知识用于读论文和复现实验}
读相关论文时，可以把本章变成一个检查表。第一，论文是否清楚定义了输入规模、模型规模、数据来源和评估协议；第二，论文声称的改进落在哪个变量上；第三，作者是否报告了会让结论失效的设置；第四，实验是否能分离模型结构、数据质量、优化设置和系统实现的影响。
复现实验时，也应避免直接追求大规模。先用最小程序复现一个公式、一个伪代码分支或一个 profile 现象，再扩大输入规模。这样做的原因很简单：如果小规模程序无法解释中间量，大规模训练只会把错误隐藏在日志和随机波动里。
把这些检查合在一起，本章给出的不是一个固定答案，而是一种阅读 CS336 的方法：每个模型机制都要落到变量，每个变量都要落到测量，每个测量都要说明协议，每个协议都要承认边界。
因此，本章中的任何一句结论都可以被改写成一个可引用的完整句式：在给定数据、模型、硬件和评估协议下，某个机制改变了某个变量，并通过某个观测指标表现出来。这样的句子比单独背诵术语更长，但它包含了自学、复现和引用时真正需要的信息。
把本章用于自学时，可以按三遍阅读。第一遍只追踪对象和定义，确认每个术语到底指向文本、token、张量、样本、请求还是评估记录。第二遍追踪公式和伪代码，确认每一步改变了哪些状态。第三遍追踪实验和 caveat，确认哪些结论只在特定规模、数据或硬件条件下成立。
把本章用于复习时，则应反过来操作：先从一个失败案例或异常指标出发，再回到相应机制。loss 曲线异常可能指向数据、优化或模型容量；吞吐异常可能指向 kernel、通信或调度；benchmark 异常可能指向 contamination、prompt protocol 或采样设置。能够从异常反推机制，是掌握本章的实用标准。
\section{例题：把概念变成可计算检查}
\subsection{70B/15T 训练时间的数量级估算}
用近似公式 `training FLOPs ≈ 6ND`，设 `N=70B`、`D=15T`、`1024` 张 GPU、单卡峰值约为 `1e15 FLOPs/s`、MFU 为 `0.4`。
\begin{itemize}
\item `6ND ≈ 6.3e24 FLOPs`，这是训练总账本。
\item 有效吞吐约为 `1024 × 1e15 × 0.4 = 4.096e17 FLOPs/s`。
\item 训练时间约 `1.54e7` 秒，折合约 `178` 天；若假设或 MFU 变化，结论会线性变化。
\end{itemize}
\textbf{结论。} napkin math 能快速暴露不可能的训练声明，也能提示优化方向是算力、并行效率还是数据规模。
这个例题的目的不是替代完整工程实现，而是给出一种最小推理形式：把问题转成变量，写出近似公式或伪代码，再说明近似何时失效。
\section{实验设计：从小例子到可复现结论}
学习 PyTorch、einops 与资源核算 不能停在概念层。一个教材化结论至少需要四件事：可构造的输入、可观测的中间量、可比较的指标，以及清楚的失败条件。下面的实验设计不是要求训练大模型，而是要求读者用小规模程序复现本章机制。
\subsection{最小输入}
先构造只触发本章核心机制的小输入。例如 tokenizer 章节可以用几个包含 rare words 和 Unicode 字符的字符串；GPU 章节可以用一个 elementwise kernel 和一个 GEMM；data 章节可以用十几篇近重复文档。小输入的价值在于它让错误可见。
\subsection{资源账本}
对同一输入写下 FLOPs、bytes moved、显存峰值、通信量或样本数量的估算。估算不需要一开始就精确，但必须说明单位和近似假设。随后用 profiler、benchmark、validation loss 或人工检查去验证估算是否同量级。
\subsection{单变量消融}
一次只改变一个因素，例如 vocabulary size、head 数、batch size、filter threshold、reward scale、decoding temperature 或 GPU 数量。若多个变量同时变化，实验只能说明系统变了，不能说明机制是什么。
\subsection{失败条件}
最后要刻意触发失败：极端 context length、错误 dtype、过小 batch、污染 benchmark、低质量数据或不可靠 verifier。失败条件常常比成功曲线更能说明一个方法的边界。
\subsection{记录与报告}
实验记录应能让另一个读者复现同一结论。最少要写清楚输入构造、代码版本、随机种子、硬件、batch shape、数据版本、评价指标和异常处理方式。报告结论时不要只写“更快”或“更好”，而要说明快在哪里、好在哪个指标上，以及是否牺牲了内存、稳定性、数据质量或可解释性。
常见报告错误是把局部现象写成普遍规律。例如一次小模型实验里的 loss 改善，不等于更大模型或不同数据分布上也会改善；一次 kernel benchmark 的吞吐提升，不等于端到端 serving latency 一定降低。教材中的实验报告应始终把结论限定在可复现条件内。
\section{课程资料与回看路径}
教材正文已经把主要概念、公式和实现逻辑组织成连续叙述；本节只提供回到课程材料的阅读路径。读者复习时可以先完成前文练习，再按下面的时间段或资料组核对细节。
\subsection{视频回看路径}
\begin{longtable}{p{0.20\linewidth}p{0.46\linewidth}p{0.26\linewidth}}
\toprule
\textbf{时间} & \textbf{关键词} & \textbf{复习任务}\\
\midrule
00:00:05-00:06:37 & number, going, talk, tensors, compute, point, loss & 回看定义、例子、推导或 caveat 在该段中的位置。\\
00:19:17-00:25:45 & going, matrix, dimensions, then, there's, hidden, think & 回看定义、例子、推导或 caveat 在该段中的位置。\\
00:38:27-00:44:55 & going, that's, flops, then, memory, number, your & 回看定义、例子、推导或 caveat 在该段中的位置。\\
00:51:02-00:57:26 & arithmetic, intensity, memory, going, your, then, maybe & 回看定义、例子、推导或 caveat 在该段中的位置。\\
01:10:50-01:17:01 & memory, store, times, number, your, that's, general & 回看定义、例子、推导或 caveat 在该段中的位置。\\
\bottomrule
\end{longtable}
这些时间段用于定位视频中的讲解顺序。严肃引用时，应回到视频片段确认讲者表述、板书或幻灯片内容。
回看视频时，不必逐字抄录字幕。更有效的方法是把每个时间段判定为定义、例子、推导、代码解释或 caveat，再把它放回本章对应小节。
\subsection{课程资料阅读路径}
\begin{longtable}{p{0.14\linewidth}p{0.48\linewidth}p{0.30\linewidth}}
\toprule
\textbf{资料组} & \textbf{标题/页块} & \textbf{关键词}\\
\midrule
official\_01 & Announcements: & https, train, resources, memory, operations, read, guide\\
official\_03 & Google Brain developed brain floating point (bf16) in 2018 to address this issue. & bf16, memory, nvidia, range, fp32, training, brain\\
official\_05 & Model FLOPs utilization (MFU) & accelerator, intensity, communication, computation, time, flop, memory\\
official\_08 & - Adam = RMSProp + momentum & memory, activations, compute, training, gradient, gradients, adam\\
official\_10 & Function: einops\_motivation & function, einops\_motivation, einops\_einsum, einops\_reduce, einops\_rearrange, tensor\_operations\_flops, arithmetic\_intensity\\
\bottomrule
\end{longtable}
课程资料通常给出更稳定的标题、公式、图或代码结构；视频则补充动机和口头解释。两者合起来构成本章的事实边界。
阅读资料时，可以把每个资料组拆成三个对象：一个定义，一个公式或伪代码动作，一个边界条件。这样比逐字翻译页面或脚本更容易形成可用知识。
\section{实践说明、历史位置与风险}
PyTorch、einops 与资源核算 在 CS336 的课程结构中连接基础机制和规模化问题。基础机制解释方法为什么成立；规模化问题检验它在更大模型、更长上下文、更复杂数据或真实 serving workload 下是否仍成立。
历史位置不等于模型年表。更重要的是识别问题如何迁移：早期方法解决小规模建模问题，现代方法常常解决同一问题在大规模数据、GPU 集群、长上下文或人类偏好约束下的新形态。
\begin{itemize}
\item 资料版本限制（version risk）：课程页、公开视频和配套资料可能在学期中更新；引用时应注明访问日期或资料版本。
\item 测量口径限制（measurement risk）：FLOPs、tokens/s、benchmark score、reward 等数字只有在协议完整时可比较。
\item 规模外推限制（scaling risk）：小模型或小 batch 的机制不必然外推到 frontier-scale。
\item 数据分布限制（data risk）：数据来源、过滤和去重策略会改变模型行为，也会改变 evaluation contamination 风险。
\end{itemize}
\section{复习要点}
下面的问题把本章内容压缩为可反复检查的条目。每个条目都应能回到前文的解释、公式、伪代码或例题。
\begin{longtable}{p{0.16\linewidth}p{0.38\linewidth}p{0.38\linewidth}}
\toprule
\textbf{条目} & \textbf{对象} & \textbf{检查动作}\\
\midrule
条目 1 & tensor（张量） & 定义它；写出一个最小例子；说明一个 failure mode；指出它影响哪个资源或评估账本。\\
条目 2 & mixed precision（混合精度） & 定义它；写出一个最小例子；说明一个 failure mode；指出它影响哪个资源或评估账本。\\
条目 3 & FLOPs（浮点运算次数） & 定义它；写出一个最小例子；说明一个 failure mode；指出它影响哪个资源或评估账本。\\
条目 4 & memory bandwidth（内存带宽） & 定义它；写出一个最小例子；说明一个 failure mode；指出它影响哪个资源或评估账本。\\
条目 5 & arithmetic intensity（算术强度） & 定义它；写出一个最小例子；说明一个 failure mode；指出它影响哪个资源或评估账本。\\
条目 6 & roofline analysis（屋顶线分析） & 定义它；写出一个最小例子；说明一个 failure mode；指出它影响哪个资源或评估账本。\\
\bottomrule
\end{longtable}
\subsection{综合自测}
\begin{itemize}
\item 我能否不用课程原句，独立解释本章中心问题？
\item 我能否把本章至少一个公式改写成可运行的估算函数？
\item 我能否指出一个看似改进但实际转移成本的方法？
\item 我能否回到视频和课程资料，找到支持本章关键论断的位置？
\item 我能否设计一个最小 ablation，验证本章一个 caveat？
\end{itemize}
\section{总结与延伸}
\subsection{本章小结}
\begin{itemize}
\item PyTorch、einops 与资源核算 的核心不是记忆术语，而是理解对象、变量和机制之间的关系。
\item 视频给出讲解顺序，课程资料提供可核对的公式、图、代码或页面结构。
\item 复习时应把每个术语连接到至少一个变量、一个实现动作和一个失败模式。
\end{itemize}
\subsection{延伸解释}
本章中的延伸解释用于连接课程材料中的概念，例如说明公式里的 proxy、解释系统 tradeoff，或把多个材料片段串成一个完整推理链。若需要严肃引用，应回到课程视频和配套资料核对原始表述。
\subsection{练习}
\begin{itemize}
\item 定义题：用中英双语解释本章任意五个重要术语，并说明它们属于 compute、memory、data、optimization、evaluation 还是 deployment 账本。
\item 推导题：选择本章一个公式，逐个解释符号、单位和近似假设，并说明哪个变量最难在真实系统中测量。
\item 实现题：把本章伪代码改写成一个最小 Python sanity check，不要求训练大模型，但要输出可检查的中间量。
\item 资料题：从“课程资料与回看路径”中选择一个视频时间段，回到原始视频核对讲者例子，并记录是否存在字幕误差。
\item 评估题：为本章方法设计一个 ablation 或 benchmark，说明控制变量、数据版本、指标和 failure criteria。
\item 边界题：说明本章一个结论在更长 context、更大 batch、更多 GPU 或更复杂数据分布下可能如何失效。
\item 综合题：把本章内容和前一章或后一章连接起来，写出一个端到端训练/推理 pipeline 中的依赖关系。
\item 批判题：指出一个容易被营销材料夸大的说法，并用本章的资源账本或评估协议反驳它。
\end{itemize}
\end{document}
